\documentclass[11pt,a4paper]{article}

\usepackage[utf8]{inputenc}
\usepackage[T1]{fontenc}
\usepackage[french]{babel}
\usepackage{lmodern}
\usepackage[margin=2cm]{geometry}
\usepackage{microtype}
\usepackage{parskip}
\usepackage{setspace}
\usepackage{titlesec}
\usepackage{xcolor}
\usepackage{enumitem}
\usepackage{tabularx}
\usepackage{array}
\usepackage{booktabs}
\usepackage{tcolorbox}
\usepackage{hyperref}

\definecolor{accent}{HTML}{D4A574}
\definecolor{inkmute}{HTML}{6A6660}
\definecolor{inksoft}{HTML}{555555}

\hypersetup{colorlinks=true, linkcolor=accent, urlcolor=accent}

\titleformat{\section}{\Large\bfseries\color{black}}{}{0pt}{}[\vspace{-0.5em}\rule{\linewidth}{0.4pt}]
\titleformat{\subsection}{\large\bfseries\color{accent}}{}{0pt}{}
\titleformat{\subsubsection}{\normalsize\bfseries}{}{0pt}{}

\setstretch{1.15}
\setlength{\parskip}{0.6em}

\title{\vspace{-2cm}\Huge\textbf{Rhétorique des Vérités Difficiles}\\[0.3em]\large\textit{Comment dire ce qui doit être dit, sans détruire celui qui doit l'entendre}}
\author{}
\date{}

\begin{document}

\maketitle
\vspace{-1.5cm}

\section*{Le constat de départ}

Avoir raison ne suffit pas. Une vérité juste, balancée sans soin, n'éclaire pas le destinataire — elle l'agresse. Et l'agression ferme les oreilles avant même que le mot soit prononcé.

Les grands rhéteurs (Cicéron, Lincoln, Churchill, Martin Luther King) ne diluent \textbf{jamais} la vérité. Ils la \textbf{présentent} de façon à ce qu'elle ne soit pas refusée.

\begin{tcolorbox}[colback=accent!10, colframe=accent, boxrule=1pt, arc=2pt]
C'est cela, l'éloquence : la science de l'emballage au service de la vérité.
\end{tcolorbox}

\section{Les 7 mécanismes rhétoriques}

\subsection{1. Posture : prescription $\rightarrow$ partage}

\textbf{Le piège :} parler en « il faut », « tu dois », « tu devrais ». Le destinataire se sent corrigé, pas accompagné.

\textbf{La technique :} convertir chaque prescription en témoignage.
\begin{itemize}[leftmargin=*, itemsep=2pt]
  \item \textit{« Ce que j'ai compris, c'est que... »}
  \item \textit{« Ce qui m'a aidé moi, c'est... »}
  \item \textit{« J'ai remarqué que... »}
  \item \textit{« Je crois sincèrement que... »}
\end{itemize}

\textbf{Effet :} la leçon devient un cadeau au lieu d'un reproche.

\subsection{2. Désamorçage explicite}

\textbf{Le piège :} sauter directement aux solutions. Le destinataire pense que tu sous-estimes la difficulté.

\textbf{La technique :} acknowledger avant d'asséner. Deux phrases suffisent :
\begin{enumerate}[leftmargin=*, itemsep=2pt]
  \item Reconnaître la difficulté.
  \item Préciser pourquoi tu le dis quand même.
\end{enumerate}

\textbf{Formule type :}
\begin{quote}
\textit{« Je sais que [X est difficile / tu es fatigué]. Je te dis ça non pas parce que je pense que c'est facile, mais parce que je crois que c'est important. »}
\end{quote}

\subsection{3. Glissement des pronoms : « tu » $\rightarrow$ « on / nous »}

\textbf{Le piège :} le \textit{tu} met en face à face. Il accuse, même quand on ne le veut pas.

\textbf{La technique :} convertir en \textit{on/nous}. « Nos pensées », « on ne contrôle pas », « il nous reste ».

\textbf{Règle bonus :} glisse au moins une fois ta propre vulnérabilité — « moi aussi je tombe dans ce piège ». Ça scelle l'inclusion.

\subsection{4. Métaphores au lieu d'instructions}

\textbf{Le piège :} énoncer une consigne abstraite.

\textbf{La technique :} remplacer l'instruction par une image.

\vspace{0.5em}
\begin{tabularx}{\linewidth}{XX}
\toprule
\textbf{Instruction directe} & \textbf{Métaphore équivalente} \\
\midrule
Arrête de comparer le présent au passé & On ne marche pas en arrière vers son avenir \\
Ton cerveau finit par le croire & Ton cerveau te le prouvera, c'est son travail \\
Sors de cette logique & C'est toi qui donnes le micro à cette voix \\
Avance pas à pas & Regarde l'horizon, pas le sillage \\
\bottomrule
\end{tabularx}

\textbf{Effet (neurosciences) :} la métaphore active \textbf{2x plus de régions cérébrales} que l'énoncé direct (Lakoff).

\subsection{5. Caution intellectuelle et historique}

\textbf{Le piège :} parler depuis ta seule autorité personnelle.

\textbf{La technique :} convoquer une figure historique qui incarne ton propos.
\begin{itemize}[leftmargin=*, itemsep=2pt]
  \item Pour la résilience : \textbf{Mandela}, 27 ans à Robben Island
  \item Pour la liberté intérieure : \textbf{Viktor Frankl}, les rescapés des camps
  \item Pour la discipline : \textbf{Marc Aurèle}, \textit{Pensées pour moi-même}
  \item Pour l'optimisme face à l'adversité : \textbf{Churchill}, \textit{We shall fight on the beaches}
  \item Pour la patience : \textbf{Léonard de Vinci}, 16 ans pour finir \textit{La Cène}
\end{itemize}

\textbf{Effet :} l'argument cesse d'être personnel pour devenir universel.

\subsection{6. Rythme : phrases courtes alternées}

\textbf{Le piège :} phrases longues en continu. L'auditeur décroche.

\textbf{La technique :} alternance délibérée. Long pour développer, court pour frapper.

\textbf{Règle du métronome.} Dans chaque paragraphe, force-toi à inclure :
\begin{itemize}[leftmargin=*, itemsep=2pt]
  \item Au moins une phrase de 3 mots
  \item Au moins une phrase de 5 mots
  \item Au moins une phrase qui dépasse 25 mots
\end{itemize}

\subsection{7. La nuance fondamentale}

\textbf{Le piège :} la binaire défaitiste. \textit{« C'est foutu »}, \textit{« personne ne m'aime »}.

\textbf{La technique :} produire la distinction conceptuelle.

\vspace{0.5em}
\begin{tabularx}{\linewidth}{XX}
\toprule
\textbf{Binaire défaitiste} & \textbf{Distinction libératrice} \\
\midrule
C'est foutu & Difficile $\neq$ foutu \\
Personne ne m'aime & Seul à cet instant $\neq$ inaimable \\
J'ai tout raté & Un échec $\neq$ une identité \\
Je n'y arriverai jamais & Pas encore $\neq$ jamais \\
\bottomrule
\end{tabularx}

\textbf{Forme finale :} \textit{« X et Y, ce n'est pas la même chose. Et toute la différence est là. »}

\newpage

\section{Les 7 exercices d'entraînement}

\begin{tcolorbox}[colback=white, colframe=accent, boxrule=0.6pt]
\textbf{Ex. 1 — Le traducteur (Posture).} 10 min/jour. Réécrire 10 phrases « il faut / tu dois » en témoignage.
\end{tcolorbox}

\begin{tcolorbox}[colback=white, colframe=accent, boxrule=0.6pt]
\textbf{Ex. 2 — Le désamorceur.} Avant chaque message difficile : 2 phrases qui acknowledgent + 1 phrase qui justifie.
\end{tcolorbox}

\begin{tcolorbox}[colback=white, colframe=accent, boxrule=0.6pt]
\textbf{Ex. 3 — Le pronom-shifter.} 10 min, 3x/semaine. Réécrire un texte en \textit{tu} en \textit{on/nous}, comparer à l'oral.
\end{tcolorbox}

\begin{tcolorbox}[colback=white, colframe=accent, boxrule=0.6pt]
\textbf{Ex. 4 — Forge de métaphores.} 20 min/semaine. Un concept abstrait, 5 métaphores, garder la meilleure. Objectif annuel : 50.
\end{tcolorbox}

\begin{tcolorbox}[colback=white, colframe=accent, boxrule=0.6pt]
\textbf{Ex. 5 — Bibliothèque des cautions.} 1/semaine. Thème + figure + anecdote + citation. Objectif annuel : 50 fiches.
\end{tcolorbox}

\begin{tcolorbox}[colback=white, colframe=accent, boxrule=0.6pt]
\textbf{Ex. 6 — Le métronome.} 5 min/jour. Réécrire un paragraphe : 1 phrase de 3 mots + 1 de 5 + 1 de 25+. Lire à voix haute.
\end{tcolorbox}

\begin{tcolorbox}[colback=white, colframe=accent, boxrule=0.6pt]
\textbf{Ex. 7 — La nuance fondamentale.} 3 min/jour. Une pensée binaire entendue, sa distinction libératrice. Cible : 10 sec de production.
\end{tcolorbox}

\subsection*{Méta-exercice hebdomadaire}

Une fois par semaine : situation difficile réelle. Version brute $\rightarrow$ appliquer les 7 mécaniques. Identifier la mécanique qui apporte le plus de gain (point fort) et celle qui résiste le plus (entraînement prioritaire).

\vspace{1em}

\begin{tcolorbox}[colback=black!90, colframe=accent, coltext=white, boxrule=1pt, arc=4pt]
\centering\large\textit{Communiquer une vérité dure avec efficacité, c'est trouver l'angle qui permet à l'autre de la recevoir sans avoir à se défendre.}
\end{tcolorbox}

\vspace{0.5em}

\textit{L'éloquence n'est pas un don. C'est une science. Et elle se travaille.}

\vspace{1em}
\hfill --- \textbf{Auguste Brain}

\end{document}
